\section{Discussion}
\label{sec:discussion}

\paragraph{Why neither FVD nor short-horizon PSNR moves at scale.}
On the standard-horizon Panda setting at 999 videos, \method changes FVD from
$154.7$ to $153.4$ and leaves PSNR essentially unchanged; both differences are
within noise. An earlier reported $5\%$ reduction is superseded by the 999-video table
and must not be reused. A $512$-parameter global modulation residual, fit on a few
conditioning frames, appears too weak a lever to shift the distribution of
generated features that FVD measures, at least at this backbone scale.

\paragraph{Why LoRA is poorly matched to single-video TTA on a 13.6B DiT.}
LoRA's adaptation surface, even at modest ranks, carries thousands to
millions of free parameters. Under a single short conditioning window---a
few seconds of video---this capacity is large enough either to overfit to
the conditioning frames within a handful of gradient steps or to fail to
align in any useful direction. Our LoRA sweeps reproduce both regimes on
\backbone, and the long-horizon Panda 1000-video evaluation shows that the
overfitting regime persists at full scale: LoRA~R8 regresses both FVD and
FID against No-TTA. \method attacks the problem from the opposite direction:
rather than searching a smaller free subspace, it reuses the frozen per-block
adaLN projections as a structured de-tying mechanism that converts a tiny
shared residual into block-specific modulation, eliminating most of the
degrees of freedom that LoRA leaves open while preserving block-level
specialization.

\paragraph{Long-horizon Panda as a controlled FVD/FID trade.}
The 1000-video long-context Panda evaluation reports
$\mathrm{FVD}\!:\!278.71\!\to\!284.14$ ($+5.43$, $+1.95\%$) for \method, with
$\mathrm{FID}\!:\!29.90\!\to\!29.53$ ($-0.36$, $-1.22\%$) and modest
improvements on PSNR, SSIM, and LPIPS. Both numbers were re-derived
independently from the per-chunk sufficient statistics; the chunked-merge
implementation agrees with a single-pass FVD/FID computation to better than
$1.5\!\times\!10^{-8}$ relative error, so this is a real property of the
generated distributions rather than a numerical artifact. The cross-method
context sharpens the interpretation: on the same setting, LoRA~R8 regresses
\emph{both} metrics ($+3.70$ FVD, $+0.42$ FID), and TinyLoRA~L24 is within
metric noise of No-TTA on both ($-0.10$ FVD, $+0.19$ FID). \method is
therefore the only TTA method here that improves per-frame Inception
statistics; its cost is concentrated in the I3D-based distributional metric.

We read this as a controlled trade between per-frame appearance and full-clip
temporal coherence. FVD's I3D backbone is sensitive to 3D spatiotemporal
structure across the entire clip, while FID's per-frame Inception backbone is
not. A test-time residual that aligns per-frame appearance to the conditioning
window can therefore improve FID and pointwise pixel metrics while displacing
the I3D distribution further from the reference. The most likely mechanism is
that the conditioning window provides a smaller fraction of the predicted
future at the 93-frame horizon, biasing the adaptation toward appearance
features and away from motion features. Designing a horizon-aware variant of
\method---for example, scheduling the adaptation update against the ratio of
conditioning frames to predicted frames, or adding a motion-coherence
regularizer to the adaptation objective---is the most natural follow-up
direction and is left for future work.

\paragraph{Limitations.}
All reported numbers are computed on 1000-video evaluation sets; smaller
discovery subsets used for hyperparameter selection are not reported in the
main results because their per-set distributional bias makes their absolute
metric values unreliable. \method has no access to future-frame supervision
and therefore relies entirely on the quality of the conditioning frames;
degenerate or very short conditioning windows could reduce its effectiveness.
We do not yet evaluate \method on text-to-video or image-to-video generation;
those are natural extensions because the same modulation pathway exists in
both formulations.
